%% notes-tex/008-xue-ffa-epistasis/sections/discussion.tex
%% [[experiments.008-xue-ffa.perspective-epistasis-in-metabolic-engineering]]
%% https://github.com/Mjvolk3/torchcell/tree/main/notes-tex/008-xue-ffa-epistasis/sections/discussion.tex
%%
%% SOURCE: hand-authored. Numbers restated here are carried from the Results and are
%% produced by the scripts named at the top of sections/results.tex.

\section{Discussion\secstatus{todo}}

The design space of a strain engineering campaign is usually described as too large to
search. Ten deletable targets already give 1{,}023 combinations, and a genome-scale target
list gives more combinations than there are strains anyone will ever build. That framing
puts the difficulty in the size of the space and implies that the remedy is throughput or a
better search.

The measurements here point at a different difficulty, which throughput does not fix. In a
space small enough to enumerate completely, with every combination and every construction
order measured, the best strain is still unreachable by the procedure that campaigns
actually use. The obstacle is not that the good combination is hard to find among many. It
is that the good combination sits behind an intermediate that a screening rule would throw
away. Of 51 improvements, 2 are reachable. Adding throughput multiplies the number of routes
a campaign can try; it does not make a route monotone.

Three consequences follow for how a campaign is designed.

\notesec{Retain losses} A rule that keeps only strains better than their parent discards the
approach to almost every improvement in this design. Carrying a fixed number of
non-improving intermediates forward per round costs build capacity in proportion, and it is
the only change to the procedure itself that addresses the reachability problem. How many to
carry is set by the depth of the valleys, which here have a median of 31\% of the base
titer.

\notesec{Build combinations, not sequences} If the intermediates are uninformative about the
endpoint, then building them is a cost with no return, and a campaign is better served by
constructing the combinations it wants directly. That is a statement about what multiplexed
editing buys: not speed at reaching a known target, but access to targets that no sequence
of single edits reaches.

\notesec{Predict the combination, not the parts} A model that predicts a combination's
phenotype from its parts' phenotypes fails on exactly the combinations that matter here,
because those are the combinations where the parts do not compose. A useful predictor has to
be trained on combinations, and the enrichment result says the existing interaction graphs
will not stand in for that training data: across 28 graph-by-model tests, the interacting
triples are never enriched for connectivity in any graph, and are significantly depleted in
five. Whatever determines which combinations interact is not recorded in physical,
regulatory, genetic, TFLink or STRING connectivity among these ten genes. Nor is it carried
by combinations measured on a different readout: a transformer trained on trigenic growth
interactions predicts neither the magnitude nor the rank of the titer interactions, and at
most their sign. The combinations a predictor learns from have to be measured on the
phenotype it is asked to predict.

\subsection*{What theory predicts about where epistasis comes from}

There is a body of theory on when deletions in a metabolic system should compose and when
they should not, and it is worth asking what it predicts here. Flux balance analysis over
the yeast metabolic network gives a trimodal distribution of pairwise epistasis and
organizes the interactions into modules that interact monochromatically, meaning that the
links between two functional modules are either all buffering or all
aggravating~\citep{segreModularEpistasisYeast2005}. A complementary result derives how
epistasis propagates in a hierarchical metabolic network with first-order kinetics, and
finds that weak epistasis at a lower level is distorted as it moves to higher-level
phenotypes, so pairwise interactions between genes should be common and should depend on
the genetic background~\citep{kryazhimskiyEmergencePropagationEpistasis2021}. Both predict
what is observed here at the level of prevalence. Interaction is common, and its sign is
consistent within a readout.

What neither predicts is the reachability result, and the reason is that both are theories
about enzymes in a flux network rather than about transcription factors upstream of one.
A factor deletion changes the expression of many pathway genes at once, and it also changes
genes outside the pathway, so the module structure a flux model assigns to the pathway is
not the structure the deletions act on. The one systematic study of epistasis between
gene-specific transcription factors in yeast reached its mechanisms through expression
profiling of the double deletions rather than through pathway
topology~\citep{sameithHighresolutionGeneExpression2015}, and proposed two routes,
buffering by induced dependency and alleviation by derepression, that connect genes in
parallel pathways with no shared reaction. The enrichment result here is consistent with
that: connectivity in a physical, regulatory, genetic or coexpression graph does not
predict which of these triples interact. Applying the metabolic theory to a factor panel
would require a model that maps a factor deletion to a distribution of flux changes first,
and that model is what the field does not yet have for titer.

\subsection*{What this design can and cannot support}

The design is complete, which is what licenses the reachability claim, and it is small, which
bounds it. Ten transcription factors in one producing background, one medium and one
harvest time do not establish a rate at which improvements are unreachable in general. They
establish that the rate can be 2 in 51 in a real production campaign's target set, which is
enough to make the assumption of composability something that has to be checked rather than
assumed.

Two features of the system may not transfer. The deleted genes are transcription factors
rather than pathway enzymes, so a deletion acts on many pathway genes at once and the
interactions may be denser than a comparable enzyme panel would give. And the phenotype is a
titer under production conditions rather than growth rate, so the deletions are under no
selection to be individually tolerable. Both cut toward more interaction here than a reader
should expect elsewhere. A single-gene-per-step campaign on enzymes may well find a monotone
route more often. What the design shows is that the question has an answer, and that the
answer is measurable rather than assumable.

The measurement itself carries one limitation that applies to every number in this
document. Each strain contributes an endpoint titer in milligrams per liter from three
replicate cultures, and the released data include no biomass measurement, so nothing here
is normalized to cell density. A deletion that lowers final biomass and a deletion that
lowers fatty acid output per cell are indistinguishable in this readout, and both appear as
a lower $f$. Two consequences follow and neither is removable by reanalysis. Interaction
scores mix specific productivity with growth, so a triple whose three deletions each cost
biomass could show a negative $\tau$ with no change in per-cell output. And the valley
depths that drive the reachability result are titer valleys, which is the right quantity
for a campaign that screens on titer, and is not evidence about the yield or the rate. What
the replicates do support is that the values are stable: the trigenic standard errors
propagate from three replicates per strain through seven strains, and the monotone-route
count moves only from 2 to 11 when a step is allowed to lose up to one standard error, so
the result is not sitting on measurement noise. A design that reported optical density
alongside titer would separate the two routes, and that is a property of the source
experiment rather than of the analysis.

\subsection*{What would settle it}

Answering the general question needs more datasets of the same shape, and the shape is
restrictive. A partial combinatorial screen cannot support a reachability claim at all,
because an unbuilt intermediate is indistinguishable from a bad one. The requirement is a
design in which every sub-combination of every reported combination was itself built and
measured, in the same background and the same batch. Screens that report only the top
performers, or only the combinations that a campaign happened to reach, do not qualify no
matter how many strains they contain.

The complementary route is theory. A production phenotype is the output of a metabolic
network under constraints, and a network model that predicts where in a pathway deletions
compose and where they do not would turn the reachability question into a calculation
performed before the strains are built. The pieces exist for enzymes in a flux
network~\citep{segreModularEpistasisYeast2005,kryazhimskiyEmergencePropagationEpistasis2021},
and what is missing is the step from a regulator deletion to the flux changes it causes.
That model does not exist yet for titer, and the enrichment result above says it will not
be recovered from connectivity alone. Building it is the direction this line of work
points.
